% SPDX-License-Identifier: Apache-2.0
\section{Ports as Named Structs}
\label{sec:x-ports}

A side of \code{run} can be a struct whose fields are the ports.
A tuple pairs its names with its types by position. Two channels of
one type can then be swapped, and the build still passes. A struct
gives each port its name where the struct declares it, where a
caller passes it and where \code{run} reads it.

\code{Diff} subtracts. Its two operands are channels of the same
type, \code{Rx<U<W>>}, declared as the fields \code{minuend} and
\code{subtrahend} of \code{DiffIn}. Its outputs are the fields of
\code{DiffOut}: a channel for the difference and a wire that is high
when the subtraction borrowed. \code{run} takes the input side apart
by field, \code{DiffIn \{ minuend, subtrahend \}}, and names the
output side \code{out}, so it sends on \code{out.diff}. Both forms
work for \code{\#[lower]}. It reads the struct where it is declared
and makes each field a port, in the order the struct declares them,
and the form says what the port is called (issue 580): a side written
\code{name: S} names its ports \code{name\_field}, so \code{out.diff}
is the port \code{out\_diff}; a side taken apart names them by the
field, so \code{minuend} is the port \code{minuend}, since there is
no side name to carry. A generic struct is read with the side's
arguments, so \code{DiffIn<W>} has ports of \code{W} bits.

\code{Top} holds a \code{Diff}. Its input side is \code{Sink}, a
role of the \code{interface!} \code{Operands}. A role is a struct as
well, and \code{\#[lower]} reads it from the \code{interface!} in the
same file. A member the role reads becomes an \code{In} or an
\code{Rx}, and a member it drives an \code{Out} or a \code{Tx}. \code{Top}
passes the operands on in a struct literal that names
\code{subtrahend} first. The netlist joins each field to the child's
port of that name, whatever order the literal uses. \code{Top} passes
its output side whole, as \code{out}, and the netlist joins its
fields in order.

A struct declared in the same file is read from the file.
\code{\#[lower]} sees one item at a time and reads the file for
anything else, as it does for the functions it inlines. A struct
declared anywhere else derives \code{Ports} instead, and the next
section shows one taken from another crate; it can only be written
\code{name: S}, since the macro cannot take apart what it cannot
see, and its ports are \code{name\_field} by the same rule. A port
name that two sides both use is refused.

The run offers a minuend every cycle and a subtrahend every other
cycle, and the sink takes nothing for three cycles. Every difference
is checked against Rust, and all six must arrive. The netlist of the
top is checked against this run under both simulators, at its ports.

\begin{figure}[htbp]
\centering
\input{ports_timing.tex}
\caption{The subtraction: the minuend on every cycle, the subtrahend
on every other, the count of differences sent, and the difference
held while the sink waits.}
\label{fig:x-ports}
\end{figure}

\lstinputlisting[caption={\code{ex\_ports.rs}. Run with \code{bazel run //lib/examples:ex\_ports}.},label={lst:x-ports}]{lib/examples/ex_ports.rs}

\lstinputlisting[style=ascii,caption={What \code{ex\_ports} printed when the build ran it: the operands offered, the count, the difference taken and the borrow, per cycle; then the lowering, two modules.},label={lst:x-ports-out}]{out_ports.txt}
